\hypertarget{evaluating-them-for-a-given-state}{%
\subsection{\texorpdfstring{Evaluating them for a given state}{Evaluating them for a given state}}\label{sec:evaluation-cost}}

The identities of Section~\ref{sec:zeta-closed-form} hold for any state, so they can be
evaluated for one. What that costs, and what it needs as input, decides whether the threshold is
a usable quantity or only a structural one.

Given the \(4^n\) Pauli expectations, \(\zeta_2\) is a single weighted sum over strings and costs
\(\Theta(4^n)\). The cubic \(\zeta_1\) runs over compatible \emph{pairs}, and per qubit exactly
ten local patterns are compatible --- \((\mathbb{I},\mathbb{I})\), \((\mathbb{I},a)\),
\((a,\mathbb{I})\) and \((a,a)\) for the three non-identity letters --- so the pair set has size
\(10^n\) and the cost is \(\Theta(10^n)\), which is \(N^{\log 10/\log 4} = N^{1.661}\) in the
input size \(N = 4^n\). Enumerating a tail of qubits vectorized while looping over a head keeps
peak memory at \(O(10^{\text{tail}})\) independently of \(n\). Measured on one core, \(\zeta_1\)
takes \(2.3\;\mathrm{s}\) at \(n = 9\), \(23.5\;\mathrm{s}\) at \(n = 10\) and
\(242\;\mathrm{s}\) at \(n = 11\), so \(n = 10\) is the largest size evaluable in under a minute;
\(\zeta_2\) is milliseconds throughout. The evaluator is exact in the sense that no sampling
enters: it reproduces \(\zeta_1 = \tfrac34 t^2 - \tfrac14 t^4\) at one qubit to \(10^{-12}\), the
maximally mixed and pure-product values to machine precision, and the Haar average of
Section~\ref{sec:zeta-closed-form} to within the sampling error of the average itself.

The input is the whole Pauli spectrum, which is the state, and no proper subset suffices --- but
the two functionals fail differently, and the asymmetry is the useful part. Their kernels pull in
opposite directions on Pauli weight. Truncating the sums to \(|P| \le w\) at \(n = 6\) on the
noisy-pure ensemble, \(\zeta_2\) is recovered to \(1.2\%\) from weight one and to \(0.3\%\) from
weight two, which is \(154\) of the \(4096\) strings, because its kernel \(14^{-|P|}\) suppresses
high weight. The \(\zeta_1\) diagonal carries \(3^{|P|}\), which amplifies it: keeping \(82\%\)
of the strings still leaves \(54\%\) error, so the top weight shell alone carries more than half
(Fig.~\ref{fig:weights}). High-weight expectations are exactly the ones local shadows estimate
worst, their variance growing as \(3^{|P|}\), so \(\zeta_1\) is the hard factor and \(\zeta_2\)
the easy one. Section~\ref{sec:pilot} shows that this does not block estimation from data,
because the estimator there never forms the spectrum at all.

\begin{figure}[tb]
\centering
\includegraphics[width=0.62\linewidth]{fig10_weight_truncation}
\caption{The two projection variances read opposite ends of the Pauli spectrum. Relative error of each functional when its sum over Pauli strings is truncated to weight $|P| \leq w$, at $n = 6$ (circles noisy-pure, squares Haar-pure). $\zeta_2$ carries the kernel $14^{-|P|}$, which suppresses high weight, so a few percent of the strings determine it. The $\zeta_1$ diagonal carries $3^{|P|}$, which amplifies high weight, so most of the spectrum is needed. The cutoff $w = n$ is exact by construction and is not plotted.}
\label{fig:weights}
\end{figure}
